\section{Software and data availability}
\label{sec:availability}

\bsi\ is released under the MIT licence. Source code, the complete test
suite, all benchmark scripts and their outputs (including every dated live
benchmark run referenced in Section~\ref{sec:live-benchmark}, retained as an
audit trail, and the validation package that produces Sections~\ref{sec:validation-study}--\ref{sec:case-study}), the documented manual-workflow protocol
(Section~\ref{sec:workflow-comparison}), the usability-study protocol
referenced in Section~\ref{sec:limitations}, and full documentation are
available at:

\begin{center}
\url{https://github.com/niloy-999/BioSeqInsight}
\end{center}

\noindent The version described in this manuscript is tagged
\texttt{v2.0.2} and archived with a permanent identifier at
\todoinline{Zenodo DOI: 10.5281/zenodo.XXXXXXX -- mint on first GitHub
Release from the v2.0.2 tag, after connecting the repository to Zenodo, and
insert the resulting DOI here before submission.} An earlier prototype
(v1.0.0), discussed for context in Sections~\ref{sec:introduction} and
\ref{sec:case-study}, is preserved unmodified at the repository tag
\texttt{v1.0-archive}.

\paragraph{Installation.}
\begin{lstlisting}
pip install bioseqinsight
\end{lstlisting}
requires Python 3.10 or later and no compiled dependencies
(Section~\ref{sec:package-architecture}). \todoinline{If a standalone
installer (Section~\ref{sec:limitations}) is available at submission time,
add its location here.}

\paragraph{Reproducing the results in this paper.}
Every quantitative result in Section~\ref{sec:evaluation} is produced by a
script included in the repository under \texttt{benchmarks/scripts/}, runnable
independently of this manuscript:

\begin{lstlisting}
git clone https://github.com/niloy-999/BioSeqInsight.git
cd BioSeqInsight && pip install -e ".[validation]"

python -m pytest --cov=bioseqinsight            # Section 5.1-5.2
python benchmarks/scripts/run_fault_injection.py         # Section 5.3
python benchmarks/scripts/run_performance_benchmark.py   # Section 5.4
python benchmarks/scripts/run_workflow_comparison.py     # Section 5.5
python benchmarks/scripts/run_structure_benchmark.py \
    --out benchmarks/results/structure_benchmark_$(date +%F)  # live service benchmark, network

# validation study (network; about a day): see validation/README.md
bash validation/run_all.sh baselines
bash validation/run_all.sh threshold
bash validation/run_all.sh casestudies
bash validation/run_all.sh report
python validation/scripts/07_manuscript_numbers.py   # every number in Sections 5.6-5.9
\end{lstlisting}

The validation study depends on the state of UniProt, RCSB and PDBe at run time
and on a frozen adversarial case file whose SHA-256 is recorded beside it; the
protocol, its amendments and the corrected analyses are in
\texttt{validation/PROTOCOL.md}. The three benchmarks not requiring network access
(Sections~\ref{sec:fault-injection}--\ref{sec:workflow-comparison}) are
deterministic and reproduce the reported figures exactly on any machine; the
live structure-resource benchmark (Section~\ref{sec:live-benchmark}) depends
on the queried resources' state at run time and is expected to vary modestly
between runs, as documented by the multiple dated runs retained in the
repository during this software's development.

\paragraph{Data availability.} No new biological data were generated. The
live benchmark candidate list
(\texttt{benchmarks/data/benchmark\_candidates.csv}) references public
UniProt accessions only; all retrieved structures originate from the public
resources cited in Section~\ref{sec:retrieval} and are subject to those
resources' own terms of use.
